\section{Extra intervals beyond the equally spaced construction (WARNING: completely AI generated and not verified formally)}
\label{sec:extra-intervals}

The first evidence for this phenomenon came from the general program
\code{circle_intervals.py}, in which all interval endpoints and lengths are
free.  Fresh runs with at most five intervals, for every
\(4\leq m\leq13\), reported the following objectives.  Every entry was
obtained with zero reported MIP gap.
\begin{center}
\begin{tabular}{@{}c cc@{}}
\toprule
\(m\) & equally spaced density & free-endpoint \(N=5\) objective \\
\midrule
4  & \(1/6\)    & \(11/64\) \\
5  & \(6/35\)   & \(6/35\) \\
6  & \(1/6\)    & \(1/6\) \\
7  & \(10/63\)  & \(10/63\) \\
8  & \(3/20\)   & \(49/320\) \\
9  & \(5/33\)   & \(136/891\) \\
10 & \(3/20\)   & \(3/20\) \\
11 & \(21/143\) & \(21/143\) \\
12 & \(1/7\)    & \(145/1008\) \\
13 & \(28/195\) & \(73/507\) \\
\bottomrule
\end{tabular}
\end{center}
The displayed objectives are rational reconstructions of the floating-point
solutions.  Thus the free-endpoint model reported a strict improvement only
for \(m=4,8,9,12,13\) in this range, and in each case its positive intervals
form the equally spaced base plus one short component.  This shape arose
without being imposed.  Lower-\(N\) cross-checks already find the improvement
with three allowed intervals for \(m=4\), four for \(m=8,9\), and five for
\(m=12,13\).  Exact rational replay verifies all five reconstructed unions as
constructions.  The artifacts are in
\code{json_results/verification_20260716/higher_m_general/} and
\code{higher_m_general_m11-13/}.  Optimality within these finite-interval
models remains a numerical result, not an exact upper bound for arbitrary
measurable sets.

\subsection{The specialized fixed-base experiment}

The script \code{equally_spaced_tiny_interval_lp.py} was written afterward to
test the shape suggested by the general program over a much larger range.  It
fixes the equally spaced base parameters \(q,\delta,a\) from
Section~\ref{sec:equally-spaced-construction}, adds a chosen number of free
intervals, and imposes the lifted constraints from
Section~\ref{sec:milp-formulation}.  The extra intervals are not assigned to
gaps in advance, so binary variables impose their ordering and disjointness.
This is a local extension model: it asks whether this fixed base can be
enlarged, not whether the base is globally optimal.

The zero-gap one-extra scan tested every \(3\leq m\leq100\), using one
representative anchor because the other anchor choices are translates.  On
the larger optimizing choice of \(q\), it found a positive gain precisely for
\(m\equiv0,1\pmod4\), except for \(m=5\).  The observed gain was
\[
 \frac1{192}\quad(m=4),\qquad
 \frac{2}{m^2(m+2)}\quad(4\mid m,\ m\geq8),
\]
and
\[
 \frac{1}{m^2(m+2)}
 \quad(m\equiv1\!\!\pmod4,\ m\geq9).
\]
The numerical gains agreed with these fractions to within
\(2\times10^{-11}\).  If \(4\mid m\), two adjacent choices of \(q\) maximize
the base density.  The smaller choice \(q=m/4\) admitted no gain above
\(10^{-9}\) in any tested case; the improvement occurs on the larger choice
\(q=m/4+1\).  All 123 saved one-extra assignments, including both tied
branches, also pass an independent replay of their bounds, integrality, and
linear constraints.

The complete two-extra scan used the larger branch for every
\(3\leq m\leq100\) and the smaller branch in every tied case, again giving
123 runs.  The first 44 were solved without a specialized upper bound.  For
the remaining 79, each extra length was bounded by the independently computed
one-extra optimum.  This is a valid redundant constraint: deleting the other
extra leaves a feasible one-extra extension of the same base.  Both extras
are positive from \(m=12\) onward in the two relevant congruence classes, and
the total gain is
\[
 \frac{4}{m^2(m+2)}\quad(4\mid m,\ m\geq12),
 \qquad
 \frac{2}{m^2(m+2)}
 \quad(m\equiv1\!\!\pmod4,\ m\geq13).
\]
For \(m=4,8,9\), the second variable has zero length.  No smaller tied branch
improves.  All 123 canonical runs report optimal status and zero stored gap;
their assignments pass independent constraint replay.  Strict no-presolve
reruns replaced noisy canonical results at \(m=34,40,100\), while the original
outputs were retained as diagnostics.  In particular, presolve at \(m=100\)
had stopped one interval exactly \(10^{-6}\) short despite reporting
optimality, illustrating why exact reconstruction and replay are necessary.

The one- and two-extra artifacts, summaries, and replay reports are in
\code{json_results/verification_20260716/tiny_intervals/} and
\code{tiny_intervals_two_extra/}.  Exact fraction arithmetic verifies every
positive two-extra output.  At \(m=12\), the two intervals have length
\(1/1008\), giving density \(73/504\); at \(m=13\), both have length
\(1/2535\), giving \(122/845\).

\subsection{More extras and exact families}

A third extra interval changes the first of these examples.  For \(m=12\),
the specialized model finds lengths
\[
 \frac1{1008},\qquad \frac5{12096},\qquad \frac1{1008},
\]
and hence the seven-interval density
\[
 \frac17+\frac{29}{12096}=\frac{251}{1728}.
\]
All 196 lifted interactions of the rationalized union have been checked
exactly.  At \(m=13\), a third variable remains zero.  At \(m=16,17\), three
extras are positive, and subsequent experiments reveal a longer regular
pattern.

To state that pattern, let \(B_m\) denote the larger optimizing equally
spaced base and put \(b=a+1/m\), with representatives chosen on the real
line.  If \(4\mid m\), set
\[
 q=\frac m4+1,
 \qquad \delta=\frac1{2(m+2)},
 \qquad g=\frac2{m^2(m+2)},
 \qquad
 s_j=b-\frac{m-4j}{4m^2}.
\]
For every tested \(m=16,20,\ldots,100\), the union of \(B_m\) with the
intervals of length \(g\) starting at
\[
 s_j\quad\hbox{and}\quad s_j+\frac1m,
 \qquad 2\leq j\leq\frac{m-4}{4},
\]
is \(m\)-sum-free.  It has \((m-8)/2\) extra intervals and density
\begin{equation}
 \frac{m^3+4m^2+8m-64}{8m^2(m+2)}
 =\mu(B_m)+\frac{m-8}{m^2(m+2)}.
 \label{eq:mod0-extra-family}
\end{equation}

If \(m\equiv1\pmod4\), set
\[
 q=\frac{m+3}{4},
 \qquad \delta=\frac{m+1}{2m(m+2)},
 \qquad g=\frac1{m^2(m+2)},
 \qquad
 s_j=b-\frac{m-(4j+1)}{4m^2}.
\]
For every tested \(m=13,17,\ldots,97\), take the intervals of length \(g\)
starting at
\[
 s_j\quad\hbox{and}\quad s_j+\frac1m,
 \qquad 2\leq j\leq\frac{m-5}{4}.
\]
Together with \(B_m\), these form an exactly checked \(m\)-sum-free union
with \((m-9)/2\) extras and density
\begin{equation}
 \frac{m^3+4m^2+7m-36}{8m^2(m+2)}
 =\mu(B_m)+\frac{m-9}{2m^2(m+2)}.
 \label{eq:mod1-extra-family}
\end{equation}
The endpoint and lift verifiers checked more than \(2.28\) million images for
these two finite families using exact rational arithmetic.  Their reports are
archived in \code{tiny_intervals_mod0_doubled_tail_exact/} and
\code{tiny_intervals_mod1_full_family_exact/}.  These checks prove the listed
lower-bound constructions; they do not prove the formulas for untested
multipliers or establish local optimality.

The aggregate gains in~\eqref{eq:mod0-extra-family}
and~\eqref{eq:mod1-extra-family} are \(O(m^{-2})\), even though the number of
extra components is linear in \(m\).  They therefore do not change the
\(1/8\) limit of the equally spaced densities.  This does not show that every
possible improvement is \(O(m^{-2})\); merely proving that the best aggregate
improvement is \(o(1)\) would suffice to preserve the limit \(1/8\).

As a numerical saturation check, the fixed-base model was solved with one
more extra variable than the number eventually used in the following cases:
\begin{center}
\begin{tabular}{@{}c c c c@{}}
\toprule
\(m\) & positive extras & exact constructed density & next tested count \\
\midrule
12 & 3 & \(251/1728\) & 5 variables, still 3 positive \\
13 & 2 & \(122/845\)  & 5 variables, still 2 positive \\
16 & 4 & \(9/64\)     & 5 variables, still 4 positive \\
17 & 4 & \(769/5491\) & 6 variables, still 4 positive \\
20 & 6 & \(303/2200\) & 7 variables, still 6 positive \\
21 & 6 & \(464/3381\) & 7 variables, still 6 positive \\
\bottomrule
\end{tabular}
\end{center}
These runs requested both zero relative and zero absolute MIP gaps.  Ordering
the exchangeable extra variables is a useful symmetry reduction: for example,
the \(m=16\), five-variable model fell from 238840 branch-and-bound nodes to
3873 without changing the optimum.  The saturation results support the exact
families above, but remain finite numerical optimality statements inside the
fixed-base model.  The case \(m=12\) is a genuine exception to the equal-length
family: its middle extra has length \(5/12096\).

The remaining theoretical questions are to prove the two displayed families
uniformly, determine whether they are locally optimal for the fixed base, and
decide whether free endpoints or entirely different constructions give a
larger aggregate improvement.  In particular, it remains open whether
\(d_m(\T)\to1/8\).
